\chapter{Introdução}

Para transformar um problema real em um modelo compatível com soluções computacionais existem algumas metodologias de idealização do fenômeno físico - onde traduzimos suas características através de modelos matemáticos - a fim de discretizá-lo e resolvê-lo, como esquematizado na Figura \ref{fig:modelos}.

\begin{figure}[hb]
\centering
\includegraphics[scale=0.7]{Imagens/Modelos}
\caption{Modelagem de problemas reais para solução numérica}
{\footnotesize Fonte: Elaborada pelo autor.}
\label{fig:modelos}
\end{figure}

As simulações computacionais fazem uso de diferentes métodos e técnicas para calcular soluções envolvendo os problemas típicos da ciência e engenharia, como: método de diferenças finitas, método de elementos finitos e método de volumes finitos. O método dos elementos finitos (MEF) é um procedimento que busca soluções aproximadas para os modelos numéricos e se aplica a uma grande variedade de problemas físicos. Sua acurácia e estabilidade estão largamente estudados, o que confere ao método uma robustez e sólida confiabilidade.

Para viabilizar a solução, a modelagem divide uma geometria grande e complexa, que é submetida a carregamentos térmicos ou mecânicos em pequenas partes – denominadas \textit{elementos} - os quais passam a representar o domínio contínuo do problema. A filosofia por trás do método é reduzir um problema grande (seu objeto de estudo) em problemas menores (os \textit{elementos}) e calcular também a relação entre eles. O conjunto dos \textit{elementos} e de seus pontos de contorno – denominados nós – é conhecido como \textit{malha}. A \textit{malha}, portanto, é composta por um número finito de \textit{elementos} de comportamento bem definido. O método de elementos finitos é um método poderoso para discretizaçãoo de geometrias complexas pois não exige esforço adicional comparado a sua utilização em geometrias regulares. 
 
A precisão do método dos elementos finitos (MEF) depende da quantidade de nós e elementos, do tamanho e dos tipos de elementos da malha. Ou seja, quanto menor for o tamanho e maior for o número dos elementos em uma determinada malha, maior a precisão nos resultados da análise.

\section{Motivaç{\~ a}o}
Veículos de competição são submetidos a grandes esforços uma vez que o seu objetivo é ser mais eficaz em uma pista de corrida. As pistas planejadas para veículos de alta performance são compostas de diversos trechos curvos, principalmente as pistas de Formula SAE e a frenagem eficiente permite maior precisão nas entradas de curva. Sendo assim o sistema de freios é constantemente acionado, o que leva a um considerável aumento de temperatura dos seus componentes. O aumento descontrolado de temperatura nos discos de freio podem resultar em fadiga, vaporização do fluido de freio, rachaduras térmicas e vibrações prejudiciais ao funcionamento do sistema, o que compromete a segurança e performance do protótipo.

\section{Objetivo}
Estudar os fenômenos energéticos envolvendo os componentes dos freios e apresentar soluções computaconais a fim de proporcionar comparativos das influências das grandezas envolvidas no projeto.

\section{Metodologia}

\section{Organizaç{\~ a}o da tese}
